Ein perfekter (idealer) Stirlingmotor erreicht den Carnot-Wirkungsgrad (FoTaBe S.~172)
\begin{align}
 \eta &= 1-\frac{T_\text{k}}{T_\text{w}}
\end{align}
mit der Temperatur des Kühlwassers $T_\text{k} = \Tk$. Aufgelöst nach der oberen Temperatur:
\begin{align}
 T_\text{w} &= \TwF \\
   &= \frac{\Tk}{1-\etSO} \\
   &= \Tw \approx \result{\TwP{0}{2}}
\end{align}
Das sind
\begin{align}
 \vartheta_\text{w} &= \twF = \tw \approx \result{\twP{0}{2}}
\end{align}
